\section{Experimental design (pre-registered)}\label{sec:design}
Everything in this section was written before any experiment, in a decision record that becomes binding when it is
tagged in the repository (\texttt{prereg-v1}) before the first timed run on a candidate instance. Later changes are
new, dated decisions; results under the original plan will still be reported.

\paragraph{Machine and footprints.} One NVIDIA DGX Spark (GB10: 10 Cortex-X925 and 10 Cortex-A725 cores, Blackwell
GPU, 121\,GiB memory shared by CPU and GPU). Each timed solver runs on one Cortex-X925 core with one thread and a hard
memory cap. The EA and cuOpt get the same partner footprint: the GPU, the same three Cortex-A725 cores and the same GPU
memory limit, 12\,GiB --- set by cuOpt's measured need on the largest unit-commitment case, since cuOpt does not run
under a much tighter limit. One GPU arm runs at a time; local language-model servers that share the GPU are stopped during
confirmatory runs, and a run whose GPU clock drops (a known fault of this machine) or that sees foreign GPU activity
voids its whole block.

\paragraph{Instances.} A fresh screening pass with SCIP alone (300\,s, two seeds) over the MIPLIB 2017 benchmark
set~\citep{gleixner2021miplib} without its infeasible instances, the \texttt{ca} and \texttt{rts\_gmlc} cases of
PGLib-UC~\citep{pglibuc2019,knueven2020uc} and 60 ML4CO item-placement instances~\citep{gasse2022ml4co} defines four
strata --- MIPLIB optimisation-hard ($P(300) \ge 0.10$), MIPLIB feasibility-hard (no feasible solution on a seed),
PGLib-UC and ML4CO --- from which 32 test and 16 tuning instances are drawn by quota, split by family so that no family
is on both sides. Hardness is measured on SCIP alone, never on cuOpt, so that selection does not decide the comparison
with cuOpt. \pending{Pool to be listed after screening.}

\paragraph{Arms.} A block is one instance and seed; the SCIP seed is the same in every arm of a block.
\begin{center}\small
\begin{tabular}{lll}
\toprule
arm & what runs & role \\
\midrule
A0 & SCIP alone, with the exchange plug-in installed and a silent partner & control (H1a) \\
A0$^*$ & SCIP alone with the best static emphasis from the tuning set & control (H5) \\
A1 & SCIP $\parallel$ SCIP (second seed), no exchange & control (H4) \\
A2 & SCIP $\parallel$ cuOpt, no exchange & control (H3) \\
A3 & SCIP $\rightleftarrows$ I-TREA II on the GPU, two-way exchange & \textbf{treatment} \\
A3r & SCIP $\parallel$ I-TREA II, no exchange (subset) & live check of H2 \\
A5 & I-TREA II alone & control (H1b), mechanism (H6) \\
A5x0 & I-TREA II alone with $x = 0$ (no Rank~0) & mechanism (H6) \\
A4, A5y0, A6 & CPU EA, $y = 0$, cuOpt alone & exploratory \\
\bottomrule
\end{tabular}
\end{center}

\paragraph{Metrics.} Primary: the portfolio primal integral $P(300)$ --- the best solution held by any process of a run
at each moment --- built only from incumbents that pass the independent exact checker, with references frozen before
the confirmatory runs. Secondary: the solver's own $P(300)$, times to first feasible and to gaps of 5\,\% and 1\,\%
(censored at $T$), $P(1{,}800)$ on a subset, the number of solver incumbents that came from the EA, and a MIPfeas-style
integral. Integrity: every incumbent of every process checked; every solution posted to the solver classified as
accepted, rejected, superseded before it was tried, or pending at $T$; crashes counted per arm.

\paragraph{Hypotheses.} Lower $P(300)$ is better.
\begin{itemize}
\item \textbf{Primary (both required):} H1a, A3 $<$ A0 (the coupling beats the solver alone), and H1b, A3 $<$ A5 (it
  beats the EA alone). An intersection--union test: each at $\alpha = 0.05$, no adjustment.
\item \textbf{Secondary (Holm):} H2, A3 $<$ the uncoupled portfolio, the pointwise best of A0 and A5 in the same block
  (optimistic for the uncoupled side; checked live by A3r on a subset); H3, A3 $<$ A2; H4, A3 $<$ A1; H5, A3 $<$
  A0$^*$.
\item \textbf{Mechanism:} H6, A5 $<$ A5x0 on the EA's own $P(300)$ --- do good infeasible solutions help?
\end{itemize}

\paragraph{Analysis.} Per instance (seeds averaged), $d_i = (P_{\text{ctrl}} - P_{\text{arm}}) /
\max(P_{\text{ctrl}}, P_{\text{arm}}, 10^{-3}) \in [-1, 1]$. A two-sided Wilcoxon signed-rank test on $d_i$, and the mean
of $d_i$ with a 95\,\% percentile bootstrap interval over instances~\citep{dicicco1996bootstrap}. A comparison is
\emph{supported} if $p < 0.05$ and the interval excludes 0; \emph{negative} only if the whole interval lies below the
non-inferiority margin $+0.05$; \emph{arm worse} if it lies below 0; otherwise \emph{inconclusive}. With 32 test
instances and two seeds (1000 and 1001, disjoint from screening and tuning seeds), the minimum detectable mean $d$ is
about 0.20--0.25. The feasibility stratum is analysed on time to first feasible. Arm order within a block follows a
Williams design; blocks run in random order with strata interleaved; crashed runs are kept and truncated, never re-run,
except for pre-registered external causes, which void and re-run the whole block.

\paragraph{Gates and tuning.} Before the confirmatory runs, the EA alone must reach a first feasible solution within
twice cuOpt's time on at least half of the tuning instances where cuOpt finds one, and the GPU local search's throughput
is reported against a single-core CPU reference; if the first gate fails, the confirmatory runs are not made and the
engineering gap is reported. The EA's configuration is then chosen on the tuning set with at most 12 configurations
$\times$ 16 instances $\times$ 120\,s, every configuration logged, and frozen.
